\begin{table}[htbp]\centering\footnotesize
\caption{El límite máximo del gasto corriente estructural con la metodología nueva (mdp)}\label{cua:gcemetodo}
\begin{tabular}{>{\raggedright\arraybackslash}p{9.6cm}r}
\toprule
Concepto & Monto \\
\midrule
\textbf{A) Gasto neto pagado CP 2025} & \textbf{8.374.179,5} \\
\quad (1) Costo financiero & 1.131.312,2 \\
\quad (2) Participaciones & 1.350.864,7 \\
\quad (3) Adefas & 28.863,6 \\
\quad (4) Pensiones y jubilaciones & 1.458.485,7 \\
\quad (5) Inversión física y financiera directa GF y ECPD & 643.527,6 \\
\quad (6) Inversión física y financiera directa ECPI & 35.936,0 \\
\quad \textbf{(7) Programas sociales universales establecidos en la CPEUM} & \textbf{766.412,4} \\
\quad \textbf{(8) Servicios personales función educación} & \textbf{648.331,6} \\
\quad \textbf{(9) Servicios personales función salud} & \textbf{522.120,5} \\
\quad \textbf{(10) Servicios personales función asuntos de orden público y de seguridad interior} & \textbf{27.635,3} \\
\textbf{B) Gasto corriente estructural pagado CP 2025} & \textbf{1.760.689,7} \\
\quad Diferimiento de pagos & 75.720,6 \\
\textbf{C) Gasto corriente estructural devengado CP 2025} & \textbf{1.836.410,3} \\
\textbf{D) Límite máximo para 2027} & \textbf{2.058.457,0} \\
\midrule \textbf{Suma de los rubros (7) a (10)} & \textbf{1.964.499,8} \\
\bottomrule
\end{tabular}
\\[2pt]\begin{minipage}{0.92\textwidth}\scriptsize Fuente: CGPE 2027, edición SHCP, p. 24, extraído del PDF a \texttt{datos/gce\_limite\_2027.csv}. El renglón A excluye, además, el gasto de las empresas públicas del Estado (nota 2/ del propio cuadro, artículo 2, fracción XXIV Bis, de la LFPRH). \textbf{Los rubros en negritas no se restaban antes}; son de la Cuenta Pública 2025, no del proyecto 2027. La D resulta de actualizar C por precios (3,8 y 4,0~\%) y crecimiento real (1,9~\% cada año).\end{minipage}
\end{table}
